\section{\LaTeX\ packages for Physics}
	
	
	\subsection{The \texttt{siunitx} package}
		The \lstinline|siunitx| package provides tools for typesetting physical quantities, numbers, and units consistently in \LaTeX. Maintained by Joseph Wright, the current version (v3) provides extensive automated formatting, parsing, rounding, and decimal-aligned table columns.
		
		To use the \lstinline|siunitx| package, we include the following line in the preamble:
\begin{lstlisting}
\usepackage{siunitx}
\end{lstlisting}
		
		The \lstinline|siunitx| package enforces the BIPM (International Bureau of Weights and Measures) standards for formatting numbers and units. For physics lab reports and dissertations, it prevents critical formatting errors such as italicizing units, misaligning decimals or using incorrect spacing between values and units.
		
		\subsubsection{Fundamental macros}
			The fundamental macros provided by \lstinline|siunitx| handles the spacing as well as the font selection. These are briefly summarised below:
			\begin{enumerate}
				\item For standalone numbers
			\end{enumerate}
	
	\subsection{The \texttt{physics} package}
		The \lstinline|physics| package is a popular \LaTeX\ package created by Sergio C. de la Barrera that simplifies and standardizes the typesetting of physics, vector calculus and linear algebra equations.
		
		To use the \lstinline|physics| package, we include the following line in the preamble:
\begin{lstlisting}
\usepackage{physics}
\end{lstlisting}